% !TEX program = lualatex
\documentclass{article}

\usepackage{indentfirst}
\usepackage{amssymb}
\usepackage{amsmath}

\title{Exercises from chapter 4}

\author{Kainã Alves Tureso - 15466391}

\begin{document}
    \maketitle

    \section*{Exercise 1}

    Suppose by absurdity that $\exists w \in \mathbb{R}^n$ such that $Aw= 0w$, that is, $\lambda = 0$ is an eigenvalue and $w$ is the respective eigenvector with $w \ne 0$. So:
    \[
        Aw = 0 \implies w^\top A w = w^\top 0 = 0 
    \]
    But $w^\top A w > 0 $. Absurd. So $A$ doesn't have a $0$ eigenvalue. Therefore $A$ is invertible

    \section*{Exercise 2}
    We can show that the coercity and continuity constants are $1$, when using the energy norm:

    \subsection*{Continuity constant}
    Using the Cauchy-Schwarz inequality:
    \[
        |a(u,v)| \le ||u||_a ||v||_a \implies N_a = 1
    \]

    \subsection*{Coercity constant}
    We just need to note hat $a(u,u) = ||u||^2_a = \alpha ||u||^2_a$. So $\alpha = 1$

    Applying this to Céa's lemma, we obtain:
    \[
        ||u - u_h||_a \le ||u- v_h||_a, \forall v_h \in V_h
    \]
\end{document}